\documentclass[11pt,a4paper]{article}

% ─── Packages ───────────────────────────────────────────────────────────────
\usepackage[margin=1in]{geometry}
\usepackage{amsmath,amssymb}
\usepackage{graphicx}
\usepackage{booktabs}
\usepackage{hyperref}
\usepackage{xcolor}
\usepackage{listings}
\usepackage{microtype}
\usepackage{setspace}
\usepackage{authblk}
\usepackage{cite}
\usepackage{tabularx}
\usepackage{multirow}
\usepackage{url}
\usepackage{enumitem}
\usepackage{caption}
\usepackage{subcaption}

\hypersetup{
  colorlinks=true,
  linkcolor=blue!70!black,
  citecolor=blue!70!black,
  urlcolor=blue!70!black,
}

\lstset{
  basicstyle=\small\ttfamily,
  breaklines=true,
  frame=single,
  backgroundcolor=\color{gray!10},
}

% ─── Title ──────────────────────────────────────────────────────────────────
\title{\textbf{Fine-Tuning Whisper for American Air Traffic Control\\
Speech Recognition: A Data-Efficient Approach with\\
a Roadmap Toward RLHF and RLVR Safety Systems}}

\author{Jeffrey Suu}
\affil{Independent Researcher}
\date{February 2025}

% ────────────────────────────────────────────────────────────────────────────
\begin{document}
\maketitle

% ─── Abstract ───────────────────────────────────────────────────────────────
\begin{abstract}
Automatic speech recognition (ASR) for air traffic control (ATC) communications
remains challenging due to domain-specific vocabulary, radio channel degradation,
and the dominance of non-American training corpora in the literature.
Existing open-source ATC ASR models, trained primarily on European and Southeast
Asian controller data, exhibit a word error rate (WER) of 30.3\% on American
ATC transmissions.
We present a data-efficient fine-tuning approach that adapts
\textit{Whisper Large v3}~\cite{radford2023whisper} to American ATC using
only 55 manually transcribed clips recorded from three major US airports
(IAH, JFK, SFO) via LiveATC.net.
Preprocessing combines a 300--3400~Hz bandpass filter, EBU~R128 loudness
normalization, and five-fold data augmentation (time stretching, pitch shifting,
noise injection, gain variation).
Full fine-tuning with conservative hyperparameters (learning rate $5 \times 10^{-6}$,
five epochs, weight decay 0.01) achieves \textbf{13.7\% WER}, a
\textbf{54.8\% relative reduction} over the baseline, using 370$\times$ less data
than the closest comparable study.
We further describe a six-layer system architecture and a 36-month roadmap for
extending this ASR foundation toward an AI ATC co-pilot via
Reinforcement Learning from Human Feedback (RLHF) and Reinforcement Learning
with Verifiable Rewards (RLVR), with formal FAA separation standards as
reward constraints.
Code and model are available at \url{https://github.com/jeffreysupergoodatcoding/liveatc-pipeline}
and \url{https://huggingface.co/jeffreysuu/whisper-atc-finetuned}.
\end{abstract}

\vspace{1em}
\noindent\textbf{Keywords:} air traffic control, automatic speech recognition,
Whisper, fine-tuning, RLHF, RLVR, aviation safety, American English ASR

% ─── 1. Introduction ────────────────────────────────────────────────────────
\section{Introduction}

Air traffic control is one of the most safety-critical communication domains on
earth. Controllers manage the separation and flow of thousands of aircraft
daily, relying almost entirely on voice radio exchanges that follow strict
phraseology standards (ICAO Doc 4444, FAA Order 7110.65).
The consequences of miscommunication are severe: the 2008 Spanair crash and
the 2023 Watsonville mid-air collision both illustrate how transcription errors
or phraseology misinterpretations can contribute to fatal incidents.

Despite this criticality, progress in ATC-specific ASR has lagged behind
general speech recognition. Two obstacles are paramount.
First, the domain vocabulary is highly specialized---callsigns, runways,
altitudes, squawk codes, and procedural phrases appear at frequencies that
confuse models trained on conversational or broadcast speech.
Second, and less often acknowledged, virtually all public ATC ASR research
focuses on non-American accents: the ATCO2 corpus is European
\cite{kocour2021atco2}, and the highest-profile recent results come from
Singapore \cite{wee2023singaporeatc}.
American English ATC---with its distinct regional accents, ATIS phraseology,
and numeric formatting conventions (e.g., ``1503'' rather than
``Fifteen Zero Three'')---is systematically underrepresented.

The controller shortage amplifies this gap.
The Federal Aviation Administration (FAA) projects a deficit of over 3,000
certified professional controllers by 2032 \cite{faashortage2023}.
Accurate, low-latency ASR is a prerequisite for any AI system that could
augment human controllers, log transmissions automatically, or surface
safety alerts.

This paper makes the following contributions:
\begin{enumerate}[noitemsep]
  \item A complete end-to-end pipeline for collecting, labeling, augmenting,
        and fine-tuning on American ATC audio (Section~\ref{sec:method}).
  \item A 54.8\% relative WER reduction on American ATC using only 55 clips,
        demonstrating that targeted fine-tuning substantially outperforms
        zero-shot application of European-trained models (Section~\ref{sec:experiments}).
  \item An analysis of why parameter-efficient fine-tuning (LoRA) is currently
        incompatible with Whisper's encoder architecture and how we circumvent
        this limitation (Section~\ref{sec:lora}).
  \item A forward-looking six-layer system architecture and research roadmap
        toward an AI ATC reasoning system using RLHF and RLVR
        (Section~\ref{sec:discussion}).
\end{enumerate}

% ─── 2. Related Work ────────────────────────────────────────────────────────
\section{Related Work}
\label{sec:related}

\subsection{ATC Automatic Speech Recognition}

Early ATC ASR systems relied on hidden Markov model (HMM) hybrids with
task-specific grammars \cite{cordero2012atcasr}.
The ATCO2 project \cite{kocour2021atco2} produced a large corpus of European
ATC audio and trained end-to-end models achieving sub-20\% WER on European
data, but the corpus is not representative of American phraseology or accents.

Pellegrini et al.\ \cite{pellegrini2021atc} demonstrated that pre-trained
speech encoders can be fine-tuned on small ATC corpora, establishing the
feasibility of the approach we pursue.
Zuluaga-Gomez et al.\ \cite{zuluagagomez2022atcosim} released ATCosim, a
simulation-based corpus, showing that synthetic data can supplement limited
real recordings.

\subsection{Whisper and Domain Adaptation}

OpenAI's Whisper \cite{radford2023whisper} is a sequence-to-sequence model
trained on 680,000 hours of weakly supervised web audio.
Its multilingual, multitask design makes it a natural candidate for domain
adaptation.
Gandhe et al.\ \cite{gandhe2023whisperadapt} showed that fine-tuning Whisper
on as few as 100 hours of domain-specific speech can halve WER, while
Shon et al.\ \cite{shon2023whisperatc} adapted Whisper to ATC using a
proprietary corpus.

\subsection{Southeast Asian ATC ASR}

Wee et al.\ \cite{wee2023singaporeatc} fine-tuned Whisper Large v2 on
approximately 20,000 Singapore Changi ATC clips, achieving roughly 8\% WER.
Their work is the closest published comparison; our dataset is 370$\times$
smaller yet achieves 13.7\% WER on a disjoint accent group (American English).

\subsection{Parameter-Efficient Fine-Tuning}

Low-Rank Adaptation (LoRA) \cite{hu2021lora} reduces trainable parameters
by inserting low-rank decomposition matrices into transformer weight matrices.
It has been applied successfully to language models \cite{dettmers2023qlora}
and some speech encoders \cite{kim2023loraspeech}, but encounters
architectural incompatibilities with Whisper (see Section~\ref{sec:lora}).

\subsection{Reinforcement Learning for Dialogue and Safety}

RLHF \cite{ouyang2022rlhf} applies preference learning to steer language model
outputs toward human-preferred behaviors.
RLVR \cite{lightman2023rlvr, guo2025deepseekr1} extends this by using
automatically verifiable reward signals, enabling scalable training without
human labelers for every sample.
To our knowledge, neither approach has been applied to ATC-specific reasoning.

% ─── 3. Methodology ─────────────────────────────────────────────────────────
\section{Methodology}
\label{sec:method}

\subsection{Data Collection}

Audio was recorded from LiveATC.net, a public streaming service that
rebroadcasts live ATC radio transmissions from airports worldwide.
We targeted three high-traffic American facilities to maximize phonetic and
procedural diversity:
\begin{itemize}[noitemsep]
  \item \textbf{KIAH} --- Houston George Bush Intercontinental (Tower, Approach, Ground)
  \item \textbf{KJFK} --- New York John F.\ Kennedy International (Tower, Ground)
  \item \textbf{KSFO} --- San Francisco International (Tower, Ground, Departure)
\end{itemize}

A Node.js pipeline (Figure~\ref{fig:pipeline}) records continuous streams in
10-minute chunks, applies silence-based segmentation (500~ms silence,
$-40$~dB threshold), and filters segments by duration
(2--20 seconds).
A Next.js web interface enables human reviewers to play each segment, verify
or edit Deepgram's draft transcription, and mark the segment as approved.
After review, 55 high-quality clips were approved.

\begin{figure}[h]
\centering
\begin{verbatim}
LiveATC.net -> Record (10-min chunks)
           -> Silence segmentation
           -> Upload to Supabase
           -> Human labeling (web UI)
           -> Export: audio + metadata.jsonl
           -> Preprocessing
           -> Fine-tuning
\end{verbatim}
\caption{End-to-end data pipeline from raw stream to fine-tuned model.}
\label{fig:pipeline}
\end{figure}

\subsection{Audio Preprocessing}
\label{sec:preprocess}

ATC transmissions are bandlimited by VHF radio hardware.
To replicate this characteristic during training and reduce the model's
sensitivity to broadband artifacts, we apply a Butterworth bandpass filter
with cutoff frequencies of 300~Hz and 3400~Hz (standard telephony band,
covering the principal energy of speech \cite{shannon1948communication}).
Following filtering, each clip is loudness-normalized to $-16$~LUFS
(EBU~R128 standard).

\subsection{Data Augmentation}
\label{sec:augmentation}

With only 55 base clips, overfitting is a primary concern.
We synthesize four additional variants per clip, yielding a 5$\times$ effective
dataset size, using the following transformations applied with randomized
parameters:

\begin{table}[h]
\centering
\caption{Augmentation parameters applied per clip.}
\begin{tabular}{ll}
\toprule
\textbf{Transform} & \textbf{Parameters} \\
\midrule
Time stretch        & Speed factor: $0.85\times$--$1.15\times$ \\
Pitch shift         & $\pm 4$ semitones \\
Gain variation      & $\pm 6$~dB \\
White noise injection & 1--8\% amplitude \\
\bottomrule
\end{tabular}
\end{table}

Crucially, the bandpass filter is re-applied \textit{after} each augmentation
step, ensuring that all training variants exhibit the same radio spectral
signature.

\subsection{Model and Training Configuration}

We fine-tune \texttt{openai/whisper-large-v3}~\cite{radford2023whisper},
the 1.55B-parameter encoder-decoder model.
Training used the HuggingFace \texttt{Seq2SeqTrainer} on a single Tesla T4
GPU (15~GB VRAM) in Google Colab.

\begin{table}[h]
\centering
\caption{Training hyperparameters.}
\begin{tabular}{ll}
\toprule
\textbf{Hyperparameter} & \textbf{Value} \\
\midrule
Base model                   & \texttt{openai/whisper-large-v3} \\
Learning rate                & $5 \times 10^{-6}$ \\
Epochs                       & 5 \\
Per-device batch size        & 4 \\
Gradient accumulation steps  & 4 (effective batch: 16) \\
Warmup steps                 & 50 \\
Weight decay                 & 0.01 \\
Precision                    & FP16 \\
Eval strategy                & Per epoch \\
Training duration            & $\approx$ 25 minutes \\
\bottomrule
\end{tabular}
\end{table}

We chose a learning rate of $5 \times 10^{-6}$---one order of magnitude below
common fine-tuning recommendations---to minimize catastrophic forgetting of
Whisper's general speech priors while encouraging American ATC specialization.
Weight decay and early stopping on validation loss serve as additional
regularizers.

% ─── 3.1 Why Not LoRA ───────────────────────────────────────────────────────
\subsection{Why LoRA Was Not Applicable}
\label{sec:lora}

Low-Rank Adaptation (LoRA)~\cite{hu2021lora} is an attractive approach for
small-data fine-tuning because it reduces the number of trainable parameters
from 1.55B to roughly 5--50M.
We attempted LoRA via the HuggingFace PEFT library~\cite{peft2023} but
encountered a fundamental incompatibility: PEFT's \texttt{get\_peft\_model}
injects adapters assuming the primary input token is \texttt{input\_ids}
(a standard language model input), whereas Whisper's encoder receives
\texttt{input\_features}---a 128-channel log-mel spectrogram tensor of shape
$(B, 128, 3000)$.
This mismatch causes a shape error at the first forward pass that is non-trivial
to patch without forking PEFT internals.

At the time of this work, no stable PEFT integration existed for
encoder-decoder speech models with Whisper's feature-extractor front-end.
We therefore proceeded with full fine-tuning, accepting the higher overfitting
risk and compensating with aggressive data augmentation and conservative
hyperparameters.

% ─── 4. Experiments ─────────────────────────────────────────────────────────
\section{Experiments}
\label{sec:experiments}

\subsection{Evaluation Protocol}

WER is computed using the \texttt{jiwer} library~\cite{jiwer2023} with
standard text normalization (lowercasing, punctuation removal).
We evaluate on a held-out split (15\% of the 55 clips, i.e., $\approx$8 clips)
not seen during training.

\subsection{Main Results}

\begin{table}[h]
\centering
\caption{WER comparison on held-out American ATC clips.}
\begin{tabular}{lcc}
\toprule
\textbf{System} & \textbf{WER (\%)} & \textbf{Rel.\ Improvement} \\
\midrule
Baseline (European ATC fine-tuned Whisper) & 30.3 & --- \\
\textbf{Ours (Whisper Large v3, full FT)}  & \textbf{13.7} & \textbf{54.8\%} \\
\bottomrule
\end{tabular}
\label{tab:mainresults}
\end{table}

The baseline is a publicly available Whisper variant fine-tuned on European
ATCO2 data.
Despite using only 55 clips, our model achieves a 16.6 percentage-point
absolute improvement.

\subsection{Qualitative Analysis}

Table~\ref{tab:qualitative} contrasts baseline and fine-tuned transcriptions
on three representative segments.

\begin{table}[h]
\centering
\caption{Qualitative transcription comparison.}
\label{tab:qualitative}
\begin{tabularx}{\linewidth}{lX}
\toprule
& \textbf{Transcription} \\
\midrule
\textbf{Reference}  & United 425 cleared for takeoff runway 28 left wind 240 at 12 \\
\textbf{Baseline}   & United four two five clear for take off runway twenty-eight left wind two forty at twelve \\
\textbf{Fine-tuned} & United 425 cleared for takeoff runway 28 left wind 240 at 12 \\
\midrule
\textbf{Reference}  & Delta 832 descend and maintain flight level 240 speed 280 knots \\
\textbf{Baseline}   & Delta eight thirty two descend maintain flight level two forty speed two eighty knots \\
\textbf{Fine-tuned} & Delta 832 descend and maintain flight level 240 speed 280 knots \\
\midrule
\textbf{Reference}  & Southwest 1503 contact Houston departure 135.7 good day \\
\textbf{Baseline}   & Southwest Fifteen Zero Three contact Houston departure one three five point seven good day \\
\textbf{Fine-tuned} & Southwest 1503 contact Houston departure 135.7 good day \\
\bottomrule
\end{tabularx}
\end{table}

The fine-tuned model consistently:
\begin{itemize}[noitemsep]
  \item Uses numeric notation for callsigns and frequencies
        (American ATC convention vs.\ spelled-out European convention)
  \item Retains the preposition ``and'' in climb/descent instructions
        (\textit{descend and maintain}, not \textit{descend maintain})
  \item Correctly formats compound headings and altimeter settings
\end{itemize}

\subsection{Comparison with Related Work}

\begin{table}[h]
\centering
\caption{Data efficiency comparison across ATC ASR systems.}
\begin{tabular}{lrrc}
\toprule
\textbf{System} & \textbf{Clips} & \textbf{WER (\%)} & \textbf{Accent} \\
\midrule
Whisper-ATC (Shon et al.)  & $\sim$2.9M utt. & 6.5 & Mixed \\
Wee et al.\ \cite{wee2023singaporeatc} & $\sim$20,000 & $\sim$8 & Singaporean \\
\textbf{Ours}              & \textbf{55}   & \textbf{13.7} & \textbf{American} \\
\bottomrule
\end{tabular}
\end{table}

Our system uses 370$\times$ fewer clips than Wee et al.\ and targets a
completely different accent region, demonstrating that targeted fine-tuning
with aggressive augmentation enables competitive performance at minimal data cost.

% ─── 5. Discussion ──────────────────────────────────────────────────────────
\section{Discussion}
\label{sec:discussion}

\subsection{Limitations}

Our evaluation set is small ($\approx$8 clips), introducing variance in WER
estimates.
The training data covers three airports; generalization to other American
facilities (e.g., Chicago O'Hare, Los Angeles, Atlanta) is untested.
Because we use full fine-tuning, catastrophic forgetting of Whisper's
multilingual and general speech capabilities may occur for out-of-domain audio.

\subsection{Toward a Full ATC AI Co-Pilot}

\subsubsection{System Architecture}

ATC transcription is only Layer~1 of a complete reasoning system.
Figure~\ref{fig:architecture} outlines a six-layer architecture that builds
on this ASR foundation:

\begin{enumerate}
  \item \textbf{Speech Recognition (this work):} Raw radio audio $\to$ text, 13.7\% WER
  \item \textbf{Natural Language Understanding:} Text $\to$ structured commands
        (airline, flight, instruction type, parameters)
  \item \textbf{World Model:} 4D airspace state tracker (position, altitude,
        speed, route, wake turbulence) integrating ADS-B, primary/secondary
        radar, TCAS, and weather radar feeds
  \item \textbf{RL Agent:} Decides optimal instructions given the world model
        state; trained first via RLHF, then refined via RLVR
  \item \textbf{Natural Language Generation:} Converts structured instructions
        to phraseology-compliant text (ICAO Doc 4444 / FAA 7110.65)
  \item \textbf{Speech Synthesis:} Converts text to radio-quality speech for
        broadcast
\end{enumerate}

\begin{figure}[h]
\centering
\begin{verbatim}
[Radio Audio]
      |
   Layer 1: ASR  (Whisper, fine-tuned)  <- this work
      |
   Layer 2: NLU  (structured command parsing)
      |
   Layer 3: World Model  (4D state: ADS-B + radar + weather)
      |
   Layer 4: RL Agent  (RLHF -> RLVR)
      |
   Layer 5: NLG  (phraseology-compliant instructions)
      |
   Layer 6: TTS  (radio-quality speech synthesis)
      |
[Broadcast / Display / Alert]
\end{verbatim}
\caption{Proposed six-layer ATC AI co-pilot architecture.}
\label{fig:architecture}
\end{figure}

\subsubsection{RLHF for ATC Reasoning}

RLHF~\cite{ouyang2022rlhf} is well-suited for Layer~4 because:
(a) expert controllers can rank candidate instructions without needing to
understand RL internals, and (b) ATC reasoning is not purely deterministic---
experienced controllers disagree on optimal sequencing strategies, making
human preference a meaningful signal.

Concretely, the pipeline would:
\begin{enumerate}[noitemsep]
  \item Present an airspace scenario (world model state) to the RL agent
  \item Generate multiple candidate instructions
  \item Have certified controllers rank instructions by safety, efficiency,
        phraseology, and situational appropriateness
  \item Train a reward model on ranked pairs
  \item Fine-tune the RL agent against the reward model via PPO~\cite{schulman2017ppo}
\end{enumerate}

\subsubsection{RLVR with Formal Safety Constraints}

For deployable safety-critical systems, verifiable rewards are preferable to
learned reward models, which can be gamed or miscalibrated~\cite{skalse2022reward}.
RLVR~\cite{lightman2023rlvr} trains models on rewards that are
automatically computable from first principles.

For ATC, candidate verifiable rewards include:
\begin{itemize}[noitemsep]
  \item \textbf{Separation compliance:} Binary/continuous reward based on
        whether projected trajectories maintain FAA minimum separation
        (3 nautical miles lateral or 1000~ft vertical in Class~B/C airspace)
  \item \textbf{Collision probability:} Estimated from physics-based
        trajectory prediction
  \item \textbf{Runway incursion risk:} Detected via surface movement model
  \item \textbf{Phraseology correctness:} Verified against a formal grammar
        derived from ICAO Doc 4444
  \item \textbf{Workload proxy:} Communication density, pending acknowledgements,
        and instruction completion latency
\end{itemize}

Advantages over RLHF: scalable (no human rater per interaction), objective,
transparent to certification authorities, and immune to reward hacking
by social persuasion.
Challenges: formal specification of ``correct'' ATC behavior covers edge cases
(emergency declarations, pilot deviations, equipment failures) that are
difficult to enumerate; multi-agent dynamics between aircraft must be simulated
accurately.

\subsubsection{Simulation Environment}

Effective RLVR training requires a high-fidelity simulator.
We propose building on BlueSky~\cite{hoekstra2016bluesky}, an open-source
Python ATC simulator, extended with:
\begin{itemize}[noitemsep]
  \item \textbf{Weather:} Wind shear, turbulence, precipitation, visibility
  \item \textbf{Sensor models:} ADS-B, primary radar, TCAS, weather radar
  \item \textbf{Aircraft performance:} Type-specific speed envelopes,
        climb/descent rates, turn radii, wake turbulence classes
  \item \textbf{Communication:} Tower, ground, approach, CTAF, ATIS,
        radio-failure protocols
\end{itemize}

\subsubsection{36-Month Development Roadmap}

\begin{table}[h]
\centering
\caption{Proposed development phases.}
\begin{tabular}{clp{7.5cm}}
\toprule
\textbf{Phase} & \textbf{Months} & \textbf{Milestones} \\
\midrule
1 & 1--6   & Scale ASR dataset to 1,000+ clips across 10+ airports;
              reduce WER below 10\% \\
2 & 7--12  & NLU parser; structured command extraction; world model prototype \\
3 & 13--18 & BlueSky integration; RLHF data collection;
              reward model training \\
4 & 19--24 & RLVR with separation and collision rewards;
              ablation studies \\
5 & 25--30 & NLG and TTS layers; end-to-end system evaluation in simulation \\
6 & 31--36 & Shadow mode testing with certified controllers;
              safety audit; publication \\
\bottomrule
\end{tabular}
\end{table}

\subsection{Open Research Questions}

\begin{enumerate}
  \item \textbf{Human--AI teaming}: How should an AI co-pilot communicate
        uncertainty to a human controller without increasing cognitive load?
  \item \textbf{Explainability}: ATC instructions must be explainable to pilots
        and investigators; current RL agents provide no rationale.
  \item \textbf{Certification}: FAA DO-178C and DO-200B apply to avionics
        software; no clear path exists for certifying learned models.
  \item \textbf{Adversarial robustness}: ATC radio is unencrypted and
        potentially susceptible to injection attacks; model robustness to
        adversarial audio is an open problem.
  \item \textbf{Multi-sector coordination}: Real-world ATC involves
        hand-offs between sectors and facilities; multi-agent RL is required.
\end{enumerate}

\subsection{Broader Impact}

\paragraph{Benefits.}
Improved ATC ASR reduces transcription errors in safety logs, enables
real-time phraseology feedback for trainee controllers, and provides a
data-efficient path toward AI-assisted ATC that could partially address
the controller shortage.

\paragraph{Risks.}
Overreliance on AI transcription or decision support without adequate human
oversight could introduce new failure modes.
ATC is an adversarial-adjacent environment (pilot errors, intentional
interference); models must be robustly validated before any operational use.
We emphasize that the system described in Section~\ref{sec:discussion} is
a research prototype; no claims of operational readiness are made.

% ─── 6. Conclusion ──────────────────────────────────────────────────────────
\section{Conclusion}

We demonstrated that Whisper Large v3 can be effectively adapted to American
ATC communications using only 55 manually transcribed clips, achieving 13.7\%
WER---a 54.8\% relative improvement over the European-trained baseline.
The key enabling factors are: (a) domain-matched preprocessing (300--3400~Hz
bandpass, EBU~R128 normalization), (b) aggressive data augmentation (5$\times$),
and (c) conservative full fine-tuning hyperparameters to minimize overfitting.

We also documented a practical incompatibility between the PEFT LoRA
implementation and Whisper's audio encoder, which may affect other researchers
attempting parameter-efficient adaptation of encoder-decoder speech models.

Looking forward, we outlined a six-layer system architecture and 36-month
roadmap for extending this ASR work toward a full AI ATC co-pilot using RLHF
and RLVR with formal FAA safety constraints.
We release all code, the fine-tuned model, and training notebooks to support
the research community.

% ─── Acknowledgements ───────────────────────────────────────────────────────
\section*{Acknowledgements}

The author thanks the operators of LiveATC.net for providing public access
to ATC audio streams, and Google for free Colab GPU compute.

% ─── References ─────────────────────────────────────────────────────────────
\bibliographystyle{plain}
\bibliography{references}

\end{document}
